% Part: normal-modal-logic
% Chapter: axioms-systems
% Section: introduction

\documentclass[../../../include/open-logic-section]{subfiles}

\begin{document}

\olfileid{nml}{axs}{int}

\olsection{ভূমিকা}

মোডাল মডেলের সাহায্যে মৌলিক মোডাল ভাষার একটি অর্থতত্ত্ব এবং !!a{formula} বৈধ হওয়ার একটি ধারণা আমাদের আছে: সব মডেলের সব জগতে সত্য হওয়া। মডেল বা কাঠামোর কোনো শ্রেণির সাপেক্ষেও বৈধতা নির্ধারণ করা যায়: সেই শ্রেণির সব মডেলের সব জগতে, অথবা নির্দিষ্ট কাঠামোর উপর প্রতিষ্ঠিত সব মডেলের সব জগতে, সত্য হওয়া। যুক্তিবিদ্যায় সাধারণত বৈধতার এই অর্থতাত্ত্বিক বর্ণনার সঙ্গে !!{derivability}-এর প্রমাণতাত্ত্বিক ধারণার সংযোগ ঘটানো হয়। আমাদের লক্ষ্য এমন একটি পদ্ধতিতে !!{derivability} সংজ্ঞায়িত করা, যাতে !!a{formula} ঠিক তখনই !!{derivable} হয়, যখন তা বৈধ।

সবচেয়ে সহজ এবং ঐতিহাসিকভাবে প্রাচীনতম !!{derivation}-পদ্ধতিগুলি হলো তথাকথিত হিলবার্ট-ধরনের বা স্বতঃসিদ্ধমূলক !!{derivation}-পদ্ধতি। বহু মোডাল যুক্তিবিদ্যার জন্য হিলবার্ট-ধরনের !!{derivation}-পদ্ধতি নির্মাণ করা তুলনামূলক সহজ; অধিতাত্ত্বিক আলোচনার বিষয় হিসেবেও এগুলি সরল (যেমন, বিশুদ্ধতা ও পূর্ণতা প্রমাণের ক্ষেত্রে)। তবে স্বাভাবিক নিষ্পাদন-পদ্ধতির তুলনায় এগুলিতে !!{formula}s প্রমাণ করা অনেক কঠিন।

হিলবার্ট-ধরনের !!{derivation}-পদ্ধতিতে !!a{formula}-এর !!a{derivation} হলো !!{formula}s-এর একটি ক্রম: কয়েকটি স্বতঃসিদ্ধ থেকে অল্প কয়েকটি অনুমান-বিধি প্রয়োগ করে ক্রমটি আলোচ্য !!{formula}-এ পৌঁছায়। আমরা যেহেতু !!{derivation}-পদ্ধতিকে অর্থতত্ত্বের সঙ্গে মেলাতে চাই, তাই নিশ্চিত করতে হবে যে !!{derivable} সূত্রগুলির সমষ্টি সব মডেলে সত্য (অথবা যেসব মডেলে সব স্বতঃসিদ্ধ সত্য, সেসব মডেলে সত্য)। এর জন্য প্রয়োজনীয় মোডাল যুক্তিবিদ্যার কিছু বৈশিষ্ট্য প্রথমে আলাদা করব: এগুলিই ‘স্বাভাবিক’ মোডাল যুক্তিবিদ্যা। স্বাভাবিক মোডাল যুক্তিবিদ্যার ক্ষেত্রে মাত্র দুটি অনুমান-বিধি ধরে নিতে হয়: মোডাস পোনেন্স এবং আবশ্যকীকরণ। স্বতঃসিদ্ধ হিসেবে নিই সর্বতঃসত্য সূত্রের সব (প্রতিস্থাপনজাত) উদাহরণ এবং আলোচ্য মোডাল যুক্তিবিদ্যা অনুযায়ী কয়েকটি মোডাল স্বতঃসিদ্ধ। সব মডেলের শ্রেণিই আলোচ্য হলেও~$\Ax{K}$\iftag{notprvDiamond}{}{ এবং~$\Ax{Dual}$}-এর সব প্রতিস্থাপনজাত উদাহরণকেও স্বতঃসিদ্ধ হিসেবে ধরতে হবে। শুধু এতেই ন্যূনতম স্বাভাবিক মোডাল যুক্তিবিদ্যা~$\Log K$ পাওয়া যায়।

\begin{defn}
\emph{মোডাস পোনেন্স}-এর বিধি হলো এই অনুমান-ছক:
\begin{prooftree}
\AxiomC{$!A$}
\AxiomC{$!A \lif !B$}
\RightLabel{\MP}
\BinaryInfC{$!B$}
\end{prooftree}
$!C \ident !A \lif !B$ হলে এবং কেবল তখনই বলি যে !!{formula}s $!A$, $!C$ থেকে মোডাস পোনেন্স দ্বারা !!a{formula}~$!B$ পাওয়া যায়।
\end{defn}

\begin{defn}
\emph{আবশ্যকীকরণ}-এর বিধি হলো এই অনুমান-ছক:
\begin{prooftree}
\AxiomC{$!A$}
\RightLabel{\Nec}
\UnaryInfC{$\Box !A$}
\end{prooftree}
$!B \ident \Box !A$ হলে এবং কেবল তখনই বলি যে !!{formula}s $!A$ থেকে আবশ্যকীকরণ দ্বারা !!{formula}~$!B$ পাওয়া যায়।
\end{defn}

\begin{defn}
স্বতঃসিদ্ধসমষ্টি~$\Sigma$ থেকে একটি \emph{!!{derivation}} হলো !!{formula}s $!B_1$, $!B_2$, \dots, $!B_n$-এর একটি ক্রম, যেখানে প্রতিটি $!B_i$ নিচের যেকোনো একটি:
\begin{enumerate}
\item কোনো সর্বতঃসত্য সূত্রের প্রতিস্থাপনজাত উদাহরণ; অথবা
\item $\Sigma$-র কোনো !!a{formula}-এর প্রতিস্থাপনজাত উদাহরণ; অথবা
\item $j$, $k < i$-বিশিষ্ট দুটি !!{formula}s $!B_j$, $!B_k$ থেকে মোডাস পোনেন্স দ্বারা প্রাপ্ত; অথবা
\item $j < i$-বিশিষ্ট !!a{formula}~$!B_j$ থেকে আবশ্যকীকরণ দ্বারা প্রাপ্ত।
\end{enumerate}
এমন কোনো !!a{derivation}-এ $!B_n \ident !A$ হলে বলি, $!A$ হলো $\Sigma$ থেকে \emph{!!{derivable}}; সংকেতে $\Sigma \Proves !A$।
\end{defn}

এই সংজ্ঞা অনুসারে !!{derivable} সূত্রগুলির সমষ্টি একটি স্বাভাবিক মোডাল যুক্তিবিদ্যা হবে এবং যে মডেলে প্রতিটি স্বতঃসিদ্ধ সত্য, সেই মডেলে প্রতিটি !!{derivable} !!{formula} সত্য হবে। !!{derivation}s-এর এই বৈশিষ্ট্যকে \emph{বিশুদ্ধতা} বলে। বিপরীত বৈশিষ্ট্য, \emph{পূর্ণতা}, প্রমাণ করা আরও কঠিন।

\end{document}
